% ============================================================================
% Generado por generar_anexo_a3.py a partir de docs/entregables/datos/.
% Los acentos se restituyeron despues de generarlo: regenerar lo revierte.
% ============================================================================
\section{Anexo. Registro del ejercicio de card sorting}

Este anexo reproduce el registro de las ocho sesiones aplicadas el 9 de agosto de 2026. Se incluye para que cualquier lector pueda revisar el material sobre el que se construyeron las conclusiones de la sección 2 sin depender de los archivos del repositorio. El contenido procede de las salidas crudas archivadas en \texttt{docs/entregables/datos/a3\_sorts/} y se reproduce sin editar; las únicas normalizaciones aplicadas fueron la restitución de los signos diacríticos y de la letra eñe, que el registro original no conservaba.

\subsection{Protocolo de condicionamiento}

Cada evaluador recibió la ficha completa de una persona de la Actividad 1 y el mazo de 35 tarjetas en orden barajado, sin acceso a la arquitectura propuesta ni a las respuestas de los demás. La plantilla de condicionamiento, común a las ocho sesiones, fue la siguiente.

\begin{uxnota}[Plantilla de condicionamiento por persona]
Eres \textit{[nombre de la persona]}, de \textit{[edad]} años, \textit{[ocupación]}. Antecedentes: \textit{[antecedentes de la ficha de A1]}. Objetivos: \textit{[objetivos]}. Puntos de dolor: \textit{[pain points]}. Hábitos con herramientas de datos: \textit{[hábitos]}. Frase característica: \textit{[cita de la persona]}.

Se te presentan 35 tarjetas de contenido de un portal centralizado de datos financieros. Agrúpalas como tú las organizarías para tu trabajo diario. Tú creas los grupos y tú les pones nombre, con el vocabulario que usarías en tu área. Cada tarjeta va en un grupo y las 35 deben quedar colocadas. Agrupa según tu trabajo y tus dolores, no según una taxonomía correcta.

Al terminar responde: qué tarjeta te costó más ubicar y por qué; qué tarjetas pondrías en dos grupos a la vez y por qué; y qué etiquetas no se entienden en tu vocabulario de trabajo. Responde en primera persona y justifica cada decisión desde tus antecedentes.
\end{uxnota}

Concluido el sorteo, y solo entonces, se presentó la estructura de cuatro módulos propuesta y se pidió resolver las cinco tareas de la prueba de árbol registradas en el apartado final de este anexo.

\subsection{Agrupaciones y vocabulario de cada evaluador}

Las tablas siguientes reproducen los grupos que construyó cada evaluador y el nombre que les dio. Los nombres se transcriben literalmente porque constituyen la evidencia de vocabulario sobre la que se apoya la sección 2: ninguno fue sugerido por el instrumento.

\subsubsection{P1. Laura Méndez (operativo)}

\begin{uxtabla}{|L{4.1cm}|Y|}{Agrupaciones de Laura Méndez (P1)}
  \uxheadrow \thd{Nombre que le dio} & \thd{Tarjetas que reunió} \\
  Buscar el dato & Buscador unificado; Catálogo de fuentes; Filtros avanzados \\
  Mis cosas guardadas & Búsquedas recientes; Favoritos; Mis alertas; Mi perfil \\
  Los temas que consulto & Cartera de crédito; Liquidez; Derivados; Captación; Mercado de dinero; Acciones \\
  Indicadores y avisos de riesgo & Tablero ejecutivo; Indicadores de riesgo; Señales de riesgo \\
  Para saber si el dato sirve & Diccionario de datos; Metadatos; Linaje; Calidad de datos; Reglas de negocio; Monitoreo de cargas \\
  Bajar y compartir el resultado & Vista previa; Exportación CSV/Excel; Compartir consulta; Historial de exportaciones \\
  Pedir acceso y permisos & Solicitar acceso; Estado de solicitudes; Flujos de aprobación; Permisos \\
  Cosas de sistemas que no me tocan & Usuarios y roles; Bitácora de accesos; Credenciales API; Consumo por API; Conectores de fuentes \\
\end{uxtabla}

\textbf{Tarjeta más difícil de ubicar.} Monitoreo de cargas suena a pantalla de sistemas y mi primer impulso fue mandarla con lo que no me toca. Pero es la respuesta a la pregunta que hago todos los días: si el dato ya está actualizado o todavía no. Terminé poniéndola con lo que uso para confiar en la cifra, aunque sigo sintiendo que quedó en un lugar raro.

\subsubsection{P2. Diego Hernández (analista)}

\begin{uxtabla}{|L{4.1cm}|Y|}{Agrupaciones de Diego Hernández (P2)}
  \uxheadrow \thd{Nombre que le dio} & \thd{Tarjetas que reunió} \\
  Buscar y descubrir fuentes & Buscador unificado; Catálogo de fuentes; Búsquedas recientes; Favoritos \\
  Entender el dato antes de usarlo & Diccionario de datos; Metadatos; Linaje; Calidad de datos; Reglas de negocio; Monitoreo de cargas \\
  Los datos del negocio & Cartera de crédito; Liquidez; Derivados; Captación; Mercado de dinero; Acciones \\
  Armar y repetir mi consulta & Filtros avanzados; Vista previa; Compartir consulta \\
  Sacar los datos: descargas y API & Exportación CSV/Excel; Historial de exportaciones; Credenciales API; Consumo por API \\
  Tableros y avisos & Tablero ejecutivo; Indicadores de riesgo; Señales de riesgo; Mis alertas \\
  Accesos: pedirlos y rastrearlos & Solicitar acceso; Estado de solicitudes; Flujos de aprobación; Permisos; Bitácora de accesos \\
  Configuración: mi cuenta y lo que mueve TI & Mi perfil; Usuarios y roles; Conectores de fuentes \\
\end{uxtabla}

\textbf{Tarjeta más difícil de ubicar.} Monitoreo de cargas. Por nombre suena a operación de TI y debería ir con conectores y administración, pero en la práctica es la primera pregunta que hago cada mañana: si la carga de ayer no corrió, mi análisis sale mal y no me entero. Terminé poniéndola con calidad y linaje porque para mí es parte de confiar en el dato, no de administrar el portal. Si el portal la esconde en un área de administrador a la que no tengo acceso, me deja ciego.

\subsubsection{P3. Roberto Valdez (propietario de datos)}

\begin{uxtabla}{|L{4.1cm}|Y|}{Agrupaciones de Roberto Valdez (P3)}
  \uxheadrow \thd{Nombre que le dio} & \thd{Tarjetas que reunió} \\
  Reglas y definiciones oficiales del dato & Catálogo de fuentes; Diccionario de datos; Metadatos; Linaje; Reglas de negocio \\
  Calidad y actualización de la información & Calidad de datos; Conectores de fuentes; Monitoreo de cargas \\
  Quién entra a mis bases (autorizaciones) & Solicitar acceso; Estado de solicitudes; Flujos de aprobación; Permisos; Bitácora de accesos \\
  Administración del portal (lo que le toca a sistemas) & Usuarios y roles; Credenciales API; Consumo por API \\
  Buscar y consultar información & Buscador unificado; Búsquedas recientes; Favoritos; Filtros avanzados; Vista previa; Tablero ejecutivo \\
  Descargas y envío de información & Exportación CSV/Excel; Compartir consulta; Historial de exportaciones \\
  Las bases de información por materia & Cartera de crédito; Liquidez; Derivados; Captación; Mercado de dinero; Acciones; Indicadores de riesgo; Señales de riesgo \\
  Mi cuenta y mis avisos & Mis alertas; Mi perfil \\
\end{uxtabla}

\textbf{Tarjeta más difícil de ubicar.} Consumo por API. Honestamente no sé bien qué es una API y estuve por dejarla con las cosas de sistemas sin pensarla más. Pero al leer la definición, volumen y detalle de peticiones automatizadas que recibe el portal, caí en cuenta de que eso significa que hay sistemas llevándose información sin que nadie me pregunte el motivo de consulta. Si de ahí sale mi cartera de crédito, eso es un acceso a mis bases y yo debería verlo y poder frenarlo. La dejé en administración porque quien la configura es TI, pero no me deja tranquilo.

\subsubsection{P4. Elena Ruiz (auditoría)}

\begin{uxtabla}{|L{4.1cm}|Y|}{Agrupaciones de Elena Ruiz (P4)}
  \uxheadrow \thd{Nombre que le dio} & \thd{Tarjetas que reunió} \\
  Rastro y evidencia del dato & Linaje; Bitácora de accesos; Calidad de datos; Monitoreo de cargas; Historial de exportaciones \\
  Accesos, permisos y autorizaciones & Permisos; Usuarios y roles; Solicitar acceso; Estado de solicitudes; Flujos de aprobación; Credenciales API \\
  Documentación oficial del dato & Catálogo de fuentes; Diccionario de datos; Metadatos; Reglas de negocio \\
  Información financiera de las áreas & Cartera de crédito; Liquidez; Derivados; Captación; Mercado de dinero; Acciones; Indicadores de riesgo; Señales de riesgo; Tablero ejecutivo \\
  Consulta y extracción de pruebas & Buscador unificado; Filtros avanzados; Vista previa; Exportación CSV/Excel; Compartir consulta; Búsquedas recientes \\
  Mi cuenta y mis avisos & Mi perfil; Favoritos; Mis alertas \\
  Conexiones con los sistemas (territorio de TI) & Conectores de fuentes; Consumo por API \\
\end{uxtabla}

\textbf{Tarjeta más difícil de ubicar.} Conectores de fuentes. No logré decidir si es una pantalla de configuración, y entonces no me toca ni entrar, o si es un inventario de las conexiones vivas, y entonces es justo lo que necesito para saber por dónde entra y sale información del banco. La definición dice configuración de la conexión, así que la mandé al bloque de sistemas, pero me quedó incómoda: si nadie me deja ver esa lista, hay un pedazo del perímetro que no puedo auditar y yo no puedo cerrar un informe con un hueco así.

\subsubsection{P5. Jorge Mendieta (ingeniería de datos)}

\begin{uxtabla}{|L{4.1cm}|Y|}{Agrupaciones de Jorge Mendieta (P5)}
  \uxheadrow \thd{Nombre que le dio} & \thd{Tarjetas que reunió} \\
  Tuberías, cargas y conexiones & Conectores de fuentes; Monitoreo de cargas; Calidad de datos; Credenciales API; Consumo por API \\
  Diccionario y estructura del dato & Catálogo de fuentes; Diccionario de datos; Metadatos; Linaje; Reglas de negocio \\
  Consultas y descargas del usuario & Buscador unificado; Filtros avanzados; Vista previa; Exportación CSV/Excel; Compartir consulta; Historial de exportaciones \\
  Tablas de negocio que alimento & Cartera de crédito; Liquidez; Derivados; Captación; Mercado de dinero; Acciones \\
  Indicadores y avisos de riesgo & Tablero ejecutivo; Indicadores de riesgo; Señales de riesgo \\
  Accesos, permisos y solicitudes & Solicitar acceso; Estado de solicitudes; Flujos de aprobación; Usuarios y roles; Permisos; Bitácora de accesos \\
  Mi espacio & Búsquedas recientes; Favoritos; Mis alertas; Mi perfil \\
\end{uxtabla}

\textbf{Tarjeta más difícil de ubicar.} Calidad de datos. Me trabé porque tiene dos vidas. Una es operativa y es mía: el archivo de anoche no pasó la validación por un campo de fecha en texto, y eso lo veo pegado al monitoreo de cargas. La otra es de gobierno: el resultado histórico de las validaciones es parte de la ficha de la fuente y ahí vive junto al diccionario, porque es lo que le dice al analista si puede confiar en esa tabla. La dejé con las cargas porque así la vivo yo, pero si el portal solo la pone ahí, el área de negocio nunca se va a enterar de que su fuente viene sucia.

\subsubsection{P6. Arturo Castañeda (directivo)}

\begin{uxtabla}{|L{4.1cm}|Y|}{Agrupaciones de Arturo Castañeda (P6)}
  \uxheadrow \thd{Nombre que le dio} & \thd{Tarjetas que reunió} \\
  Lo que reviso antes de la Junta & Tablero ejecutivo; Indicadores de riesgo; Señales de riesgo; Mis alertas; Calidad de datos \\
  Los temas del sistema financiero & Cartera de crédito; Liquidez; Derivados; Captación; Mercado de dinero; Acciones \\
  Respaldo de la cifra & Linaje; Diccionario de datos; Metadatos; Reglas de negocio; Monitoreo de cargas \\
  Cómo encuentro algo & Buscador unificado; Búsquedas recientes; Favoritos; Filtros avanzados; Catálogo de fuentes \\
  Lo que le pido y le mando a mi equipo & Exportación CSV/Excel; Vista previa; Historial de exportaciones; Compartir consulta \\
  Accesos y autorizaciones & Solicitar acceso; Estado de solicitudes; Flujos de aprobación; Permisos; Usuarios y roles; Bitácora de accesos \\
  Cosas de sistemas que no me tocan & Credenciales API; Consumo por API; Conectores de fuentes \\
  Mi cuenta & Mi perfil \\
\end{uxtabla}

\textbf{Tarjeta más difícil de ubicar.} Monitoreo de cargas. La etiqueta suena a tablero de sistemas y mi primer impulso fue mandarla a la pila de cosas que no me tocan. Pero al leer la definición me di cuenta de que responde a mi pregunta de todos los lunes: el dato que estoy viendo ya está actualizado o me van a decir en la Junta que es de hace tres semanas. Terminé poniéndola en Respaldo de la cifra, pero dudé bastante y no me gusta cómo se llama.

\subsubsection{P7. Mariana Ovalle Ríos (administración)}

\begin{uxtabla}{|L{4.1cm}|Y|}{Agrupaciones de Mariana Ovalle Ríos (P7)}
  \uxheadrow \thd{Nombre que le dio} & \thd{Tarjetas que reunió} \\
  Cuentas y permisos & Usuarios y roles; Permisos; Mi perfil \\
  Solicitudes y autorizaciones & Solicitar acceso; Estado de solicitudes; Flujos de aprobación \\
  Evidencia para auditoría & Bitácora de accesos; Historial de exportaciones \\
  Continuidad del servicio & Monitoreo de cargas; Calidad de datos; Mis alertas; Señales de riesgo \\
  Sistemas conectados (accesos sin persona) & Conectores de fuentes; Credenciales API; Consumo por API \\
  Ficha de cada fuente & Catálogo de fuentes; Diccionario de datos; Metadatos; Linaje; Reglas de negocio \\
  Herramientas de consulta de los usuarios & Buscador unificado; Búsquedas recientes; Favoritos; Filtros avanzados; Vista previa; Exportación CSV/Excel; Compartir consulta \\
  Información de las áreas de negocio & Cartera de crédito; Liquidez; Derivados; Captación; Mercado de dinero; Acciones; Tablero ejecutivo; Indicadores de riesgo \\
\end{uxtabla}

\textbf{Tarjeta más difícil de ubicar.} Linaje. Al leer la etiqueta pensé de inmediato en cadena de custodia, que sí es mi terreno: de dónde salió, quién lo tocó y quién lo entregó. Iba directo a mi grupo de evidencia. Al leer la definición vi que habla del recorrido del dato hasta el reporte, no del recorrido de las personas, y terminé moviéndola a la ficha de la fuente. Me quedó la duda: si un día auditoría me pregunta de dónde salió una cifra que alguien exportó, no sé si la respuesta está en Linaje o en mi bitácora.

\subsubsection{P8. Ximena Solís Barrera (integración)}

\begin{uxtabla}{|L{4.1cm}|Y|}{Agrupaciones de Ximena Solís Barrera (P8)}
  \uxheadrow \thd{Nombre que le dio} & \thd{Tarjetas que reunió} \\
  El contrato del dato: qué es y de dónde sale & Catálogo de fuentes; Diccionario de datos; Metadatos; Linaje; Reglas de negocio \\
  Armar y probar la consulta antes de programarla & Buscador unificado; Filtros avanzados; Vista previa; Compartir consulta \\
  Los conjuntos que me piden del negocio & Cartera de crédito; Liquidez; Derivados; Captación; Mercado de dinero; Acciones; Tablero ejecutivo; Indicadores de riesgo \\
  Cómo sale el dato: API y descargas & Consumo por API; Conectores de fuentes; Exportación CSV/Excel; Historial de exportaciones \\
  Que no se rompa: monitoreo y avisos & Monitoreo de cargas; Calidad de datos; Mis alertas; Señales de riesgo \\
  Trámites de acceso y permisos & Solicitar acceso; Estado de solicitudes; Flujos de aprobación; Usuarios y roles; Permisos; Bitácora de accesos \\
  Mi cuenta y mis llaves & Mi perfil; Credenciales API; Favoritos; Búsquedas recientes \\
\end{uxtabla}

\textbf{Tarjeta más difícil de ubicar.} Conectores de fuentes. No sé si es MI conexión al portal o la conexión del portal con los sistemas de origen. Si es lo primero, va conmigo en integración; si es lo segundo, es de la plataforma y ni debería verla. La dejé en salida del dato por descarte, pero la definición dice sistema de origen y eso apunta a que no es mía. Es exactamente el tipo de etiqueta ambigua que me hace abrir la pantalla nada más para descubrir que no tengo permiso.

\subsection{Peticiones de etiquetado múltiple}

Se registraron 49 peticiones de colocar una misma tarjeta en dos o más grupos. La tabla reúne las tarjetas solicitadas por tres o más evaluadores, que son las que originaron las facetas transversales de la sección 3, con la razón textual de uno de los evaluadores que la pidió.

\begin{uxtabla}{|L{2.9cm}|L{1.1cm}|Y|}{Tarjetas con petición de etiquetado múltiple (3 o más evaluadores)}
  \uxheadrow \thd{Tarjeta} & \thd{Evaluadores} & \thd{Razón registrada} \\
  Consumo por API & 6 de 8 & Mi consumo de API es mi evidencia de que la extracción automatizada corrió; el consumo total del portal es control de TI. Es la misma pantalla vista desde dos permisos distintos. (Diego) \\
  Calidad de datos & 6 de 8 & El resultado de las validaciones es parte del sello de la fuente. Cuando yo digo que esta es la versión certificada por mi subgerencia, esa afirmación se sostiene en las validaciones, así que debe leerse junto a la definición y no en otra pantalla. (Roberto) \\
  Vista previa & 4 de 8 & La vista previa me sirve en dos momentos: para confirmar que la consulta trajo lo que pedí y para revisar antes de descargar. Si solo vive en la parte de descargas, la voy a extrañar mientras todavía estoy buscando. (Laura) \\
  Catálogo de fuentes & 4 de 8 & Para mí el catálogo es mi inventario, lo que gobierno. Para el analista es por donde empieza a buscar. Es la misma lista vista desde dos escritorios. (Roberto) \\
  Mis alertas & 4 de 8 & Como preferencia personal va en mi cuenta, pero lo que yo quiero que me avise es exactamente lo del grupo de calidad: cargas y valores fuera de rango. (Roberto) \\
  Credenciales API & 4 de 8 & La llave la crea y la rota sistemas, eso lo entiendo. Pero una credencial es un privilegio otorgado, y si no aparece en la misma vista donde reviso usuarios y roles se me van a escapar las llaves huérfanas de gente que ya no trabaja aquí. Es el hallazgo clásico que encuentro cada año. (Elena) \\
  Monitoreo de cargas & 3 de 8 & El proceso de carga es de sistemas, pero el resultado (si ya cargó y a qué hora) es información que yo necesito antes de usar la cifra. Me bastaría con ver la fecha y hora de actualización junto al dato. (Laura) \\
  Permisos & 3 de 8 & Permisos tiene dos caras: la del administrador que los asigna y la mía, que solo quiere entender por qué no puedo abrir una fuente y a quién se la pido. (Laura) \\
  Historial de exportaciones & 3 de 8 & El historial de exportaciones es la lista de quién se llevó una copia de mis datos. Para el analista es el seguimiento de su descarga; para mí es control de accesos puro. (Roberto) \\
\end{uxtabla}

\subsection{Etiquetas señaladas como poco claras}

Los evaluadores marcaron 48 veces alguna etiqueta como ajena a su vocabulario. La tabla ordena las más señaladas y recoge un comentario textual de cada una; son las candidatas a renombrarse antes del prototipado.

\begin{uxtabla}{|L{2.9cm}|L{1.1cm}|Y|}{Etiquetas con más señalamientos de falta de claridad}
  \uxheadrow \thd{Etiqueta} & \thd{Evaluadores} & \thd{Comentario registrado} \\
  Metadatos & 8 de 8 & Metadatos suena técnico y no distingo bien en qué se diferencia del diccionario de datos. Para mí los dos responden lo mismo: qué significa el campo y desde cuándo es válido. (Laura) \\
  Reglas de negocio & 7 de 8 & Reglas de negocio puede querer decir políticas de la institución o la fórmula con la que se calcula la cifra. Con la pura etiqueta no sé cuál de las dos es, y esa duda es justo la que me hace tomar la variable equivocada. (Laura) \\
  Señales de riesgo & 7 de 8 & Señales de riesgo contra Indicadores de riesgo contra Mis alertas: tres etiquetas casi iguales. En mi área le diríamos alertas a todo. Espero acabar recibiendo el mismo aviso por tres caminos. (Diego) \\
  Linaje & 6 de 8 & Linaje no es una palabra que yo use. Si no leo la definición no sé qué es; yo lo llamaría 'de dónde viene el dato' o 'quién lo alimenta'. (Laura) \\
  Consumo por API & 5 de 8 & Consumo por API no me dice nada. Entiendo 'consumo' como gasto y me confunde; no sabría si tiene que ver conmigo. (Laura) \\
  Credenciales API & 3 de 8 & Credenciales API me suena a contraseñas y por un momento pensé que era donde cambio la mía. Tuve que leer la definición para descartarlo. (Laura) \\
  Conectores de fuentes & 3 de 8 & Conectores de fuentes lo asocio con cables o configuración de sistemas; no sabría si me sirve para ver qué fuentes hay disponibles. (Laura) \\
  Flujos de aprobación & 3 de 8 & Flujos de aprobación no dice aprobación de qué. Puede ser de mi acceso, de una descarga pesada o de la publicación de una fuente nueva. Lo mandé al grupo de accesos porque es el único trámite que yo vivo. (Diego) \\
  Mercado de dinero & 2 de 8 & Mercado de dinero y Captación las entiendo del negocio, pero como etiquetas de un portal no me dicen la granularidad ni el periodo que traen. El nombre del silo no me sirve si no veo al lado la periodicidad y la vigencia. (Diego) \\
  Indicadores de riesgo & 2 de 8 & Indicadores de riesgo. Para mí riesgo significa riesgo operativo, de control y de acceso; aquí resulta que es exposición financiera. Es un choque de vocabulario que me hizo dudar dónde ponerla y dónde la buscaría después. (Elena) \\
\end{uxtabla}

\subsection{Registro de la prueba de árbol}

Las 40 observaciones que sostienen el apartado de pruebas de la sección 3. La columna de coincidencia indica si el primer nivel elegido corresponde al previsto por la arquitectura puesta a prueba.

\begin{uxtablalarga}{|L{1.0cm}|L{3.3cm}|L{2.3cm}|Y|L{1.4cm}|}{Registro completo de la prueba de árbol}{\thd{Ev.} & \thd{Tarea} & \thd{Primer nivel} & \thd{Segundo nivel indicado} & \thd{Coincide}}
  P1 & Localizar una fuente de liquidez & Inicio (titubeo) & Buscador, escribiendo 'liquidez' & No \\
  P1 & Exportar la cartera de crédito & Exploración & Cartera de crédito y ahí el botón de descargar a Excel & Sí \\
  P1 & Consultar un indicador de riesgo & Exploración (titubeo) & Indicadores de riesgo & No \\
  P1 & Revisar quién autorizó un acceso & Administración (titubeo) & Bitácora de accesos o flujos de aprobación & Sí \\
  P1 & Crear una credencial de integración & Administración & Credenciales API & Sí \\
  P2 & Localizar una fuente de liquidez & Exploración & Catálogo de fuentes (o el buscador unificado) & Sí \\
  P2 & Exportar la cartera de crédito & Exploración (titubeo) & Cartera de crédito, con el botón de exportar dentro del resultado & Sí \\
  P2 & Consultar un indicador de riesgo & Exploración (titubeo) & Indicadores de riesgo & No \\
  P2 & Revisar quién autorizó un acceso & Administración (titubeo) & Bitácora de accesos / Flujos de aprobación & Sí \\
  P2 & Crear una credencial de integración & Inicio (titubeo) & Mi perfil / Mi cuenta, en una pestaña de credenciales o llaves & No \\
  P3 & Localizar una fuente de liquidez & Exploración (titubeo) & Catálogo de fuentes & Sí \\
  P3 & Exportar la cartera de crédito & Exploración & Cartera de crédito, y ahí mismo la opción de descargar & Sí \\
  P3 & Consultar un indicador de riesgo & Exploración (titubeo) & Indicadores de riesgo & No \\
  P3 & Revisar quién autorizó un acceso & Gobierno (titubeo) & Bitácora de accesos o flujos de aprobación & No \\
  P3 & Crear una credencial de integración & Administración & Credenciales API & Sí \\
  P4 & Localizar una fuente de liquidez & Exploración (titubeo) & Catálogo de fuentes & Sí \\
  P4 & Exportar la cartera de crédito & Exploración & Cartera de crédito, y de ahí el botón de exportar sobre el resultado & Sí \\
  P4 & Consultar un indicador de riesgo & Gobierno (titubeo) & Indicadores y señales de riesgo & Sí \\
  P4 & Revisar quién autorizó un acceso & Gobierno (titubeo) & Bitácora de accesos, o flujos de aprobación & No \\
  P4 & Crear una credencial de integración & Administración & Credenciales API & Sí \\
  P5 & Localizar una fuente de liquidez & Exploración & Catálogo de fuentes & Sí \\
  P5 & Exportar la cartera de crédito & Exploración & Cartera de crédito, y dentro del resultado la acción de exportar & Sí \\
  P5 & Consultar un indicador de riesgo & Exploración (titubeo) & Indicadores de riesgo & No \\
  P5 & Revisar quién autorizó un acceso & Administración (titubeo) & Bitácora de accesos, o flujos de aprobación & Sí \\
  P5 & Crear una credencial de integración & Administración & Credenciales API & Sí \\
  P6 & Localizar una fuente de liquidez & Inicio (titubeo) & Buscador unificado, la caja de búsqueda de la pantalla principal & No \\
  P6 & Exportar la cartera de crédito & Exploración & Cartera de crédito, y dentro de la vista el botón de descargar & Sí \\
  P6 & Consultar un indicador de riesgo & Inicio & Tablero ejecutivo, indicadores clave del mes & No \\
  P6 & Revisar quién autorizó un acceso & Administración (titubeo) & Bitácora de accesos, o Flujos de aprobación & Sí \\
  P6 & Crear una credencial de integración & Administración (titubeo) & Credenciales API, o algo que diga sistemas & Sí \\
  P7 & Localizar una fuente de liquidez & Exploración (titubeo) & Buscador unificado o Catálogo de fuentes & Sí \\
  P7 & Exportar la cartera de crédito & Exploración & Cartera de crédito, con la opción de exportar dentro del resultado & Sí \\
  P7 & Consultar un indicador de riesgo & Exploración (titubeo) & Indicadores de riesgo & No \\
  P7 & Revisar quién autorizó un acceso & Administración (titubeo) & Bitácora de accesos, o el historial dentro de Usuarios y roles & Sí \\
  P7 & Crear una credencial de integración & Administración & Credenciales API & Sí \\
  P8 & Localizar una fuente de liquidez & Exploración & Catálogo de fuentes (o el buscador unificado desde ahí mismo) & Sí \\
  P8 & Exportar la cartera de crédito & Exploración (titubeo) & Cartera de crédito, y desde el resultado de la consulta el botón de exportar & Sí \\
  P8 & Consultar un indicador de riesgo & Exploración & Indicadores de riesgo & No \\
  P8 & Revisar quién autorizó un acceso & Administración (titubeo) & Bitácora de accesos, o el estado de la solicitud si es una que yo pedí & Sí \\
  P8 & Crear una credencial de integración & Administración (titubeo) & Credenciales API & Sí \\
\end{uxtablalarga}

